%! Sinusoidal Function Transformations — Unit 3, Lesson 3.6 — Lesson Plan
\documentclass[10pt]{article}
\usepackage{precalculus-article}
\usepackage{precalculus-boxes}
\usepackage{graphicx}
\graphicspath{{images/}}
\newcommand{\TallMath}[1]{$\displaystyle #1\rule[-1.4em]{0pt}{3.2em}$}
\newcommand{\CourseName}{Precalculus}
\newcommand{\SchoolYear}{2026--2027}
\newcommand{\MeetingLength}{55 minutes}
\newcommand{\UnitNumberName}{Unit 3: Trigonometric and Polar Functions \quad}
\newcommand{\LessonNumberName}{Lesson 3.6: Sinusoidal Function Transformations}

\begin{document}

%----------------------------------------------------------------------
% Title Block
%----------------------------------------------------------------------
\begin{center}
  {\Large\bfseries \CourseName: \SchoolYear} \\[4pt]
  {\small\bfseries \UnitNumberName \LessonNumberName \quad (2--3 instructional periods, \MeetingLength\ each)}
\end{center}

%----------------------------------------------------------------------
% Primary Objective
%----------------------------------------------------------------------
\begin{tcolorbox}[colback=sky, colframe=navy, boxrule=0.9pt, arc=2mm,
    left=2mm, right=2mm, top=1.5mm, bottom=1.5mm]
  \textbf{Primary Objective:} Students will identify the amplitude, period,
  phase shift, and midline of sinusoidal functions of the form
  $f(\theta)=a\sin(b(\theta+c))+d$ and $g(\theta)=a\cos(b(\theta+c))+d$,
  and connect each parameter to an additive or multiplicative transformation
  of the parent sine or cosine function.
  \textit{(Big Idea: Functions)}
\end{tcolorbox}

\vspace{0.08in}

%----------------------------------------------------------------------
% Priority Ideas & Skills
%----------------------------------------------------------------------
\begin{skillbox}[Priority Ideas \& Skills]{goldbox}
\begin{tabularx}{\linewidth}{@{} p{0.46\linewidth} X @{}}
\textbf{AP Skills} & \textbf{Key Understandings (EKs)} \\
\midrule
\textbf{1.C} \textit{Construct new functions using transformations} ---
  Apply amplitude, period, midline, and phase-shift transformations
  to build $f(\theta)=a\sin(b(\theta+c))+d$ from the parent sine function.
&
\textbf{EK 3.6.A.1}: $a\sin(b(\theta+c))+d$ and $a\cos(b(\theta+c))+d$
(with $a\neq 0$) are sinusoidal transformations of sine/cosine.
Because $\cos\theta=\sin(\theta+\pi/2)$, the same transformations apply
to both. \\[4pt]

\textbf{2.B} \textit{Construct equivalent representations} ---
  Translate between parameter form $a\sin(b(\theta+c))+d$ and feature
  descriptions (amplitude, period, midline, phase shift), and rewrite
  a cosine function as a phase-shifted sine function.
&
\textbf{EK 3.6.A.2}: $\sin(\theta)+d$ shifts the midline to $y=d$
(vertical translation). \\[4pt]

&
\textbf{EK 3.6.A.3}: $\sin(\theta+c)$ shifts the graph $-c$ units
horizontally (phase shift). \\[4pt]

&
\textbf{EK 3.6.A.4}: $a\sin(\theta)$ scales amplitude by $|a|$
(vertical dilation). \\[4pt]

&
\textbf{EK 3.6.A.5}: $\sin(b\theta)$ scales period by $\tfrac{1}{|b|}$;
new period $=\dfrac{2\pi}{|b|}$ (horizontal dilation). \\[4pt]

&
\textbf{EK 3.6.A.6}: $f(\theta)=a\sin(b(\theta+c))+d$ has amplitude
$|a|$, period $\dfrac{2\pi}{|b|}$, midline $y=d$, phase shift $-c$. \\
\end{tabularx}
\end{skillbox}

\vspace{0.08in}

%----------------------------------------------------------------------
% Vocabulary
%----------------------------------------------------------------------
\begin{skillbox}[{Vocabulary, Concepts, \& Theorems}]{greenbox}
\begin{tabularx}{\linewidth}{@{} p{0.24\linewidth} X @{}}
\textbf{Term} & \textbf{Definition / Note} \\
\midrule
sinusoidal function &
  Any function equivalent to $a\sin(b(\theta+c))+d$ (or cosine form)
  with $a\neq 0$. A transformation of the parent sine or cosine. \\[4pt]
amplitude &
  $|a|$ --- half the vertical distance from midline to max (or min).
  Controls vertical stretch/compression. \\[4pt]
midline &
  The horizontal line $y=d$; the average of max and min output values.
  Determined by the vertical shift $d$. \\[4pt]
period &
  Length of one complete cycle; $T=\dfrac{2\pi}{|b|}$.
  Controlled by the horizontal dilation factor $b$. \\[4pt]
phase shift &
  Horizontal translation $= -c$ when written in factored form
  $b(\theta+c)$. Positive $c$ shifts left; negative $c$ shifts right. \\[4pt]
vertical dilation &
  Multiplication of all outputs by $a$; changes amplitude and can
  reflect over the midline when $a<0$. \\[4pt]
horizontal dilation &
  Replacement of $\theta$ by $b\theta$; changes period.
  $|b|>1$ compresses; $|b|<1$ stretches. \\
\end{tabularx}
\end{skillbox}

\vspace{0.08in}

%----------------------------------------------------------------------
% Activate Prior Knowledge
%----------------------------------------------------------------------
\begin{skillbox}[Activate Prior Knowledge \& Spiral Review]{sky}
\begin{tabularx}{\linewidth}{@{}X c@{}}
\begin{minipage}[t]{0.96\linewidth}\vspace{0pt}
\textbf{Spiral topics activated during Warm-Up (first 8--10 minutes):}
\begin{itemize}[leftmargin=1.4em, itemsep=2pt]
  \item \textbf{Unit-circle exact values}: Evaluating $\sin\theta$ and
    $\cos\theta$ at $0,\,\pi/6,\,\pi/4,\,\pi/3,\,\pi/2$ and their
    equivalents --- fluency needed to verify transformations numerically.
  \item \textbf{Graphs of $y=\sin\theta$ and $y=\cos\theta$}: Students
    must recall amplitude $= 1$, period $= 2\pi$, midline $y = 0$ from
    Lessons 3.4--3.5. Today's lesson stretches and shifts those baseline graphs.
  \item \textbf{Function transformation vocabulary (Units 1--2)}: Vertical
    shifts, vertical stretches, and horizontal shifts of polynomials and
    exponentials --- the same language and algebraic structure apply here.
\end{itemize}
\end{minipage}
&
\end{tabularx}
\end{skillbox}

\vspace{0.08in}

%----------------------------------------------------------------------
% Hook
%----------------------------------------------------------------------
\begin{skillbox}[Hook --- Entry Event]{sky}
\textbf{Scenario:} Project two sinusoidal graphs on the board side-by-side.
Graph~1: $y=\sin\theta$ (amplitude 1, period $2\pi$, midline $y=0$).
Graph~2: $y=3\sin(2\theta)+1$ (taller, faster, shifted up).
Both look ``wavy'' --- but they are clearly different.

\textbf{Discussion prompts:}
\begin{enumerate}[leftmargin=1.4em, itemsep=3pt]
  \item \textit{``The two graphs are related. What changed? What stayed the same?''}
  \item \textit{``In $y = 3\sin(2\theta)+1$, which number makes the wave
    taller? Which makes it cycle faster? Which shifts it up?''}
  \item \textit{``How would you shift the wave to the \emph{right} without
    changing its height or speed?''}
\end{enumerate}

\textbf{Key tension to surface:} All sinusoidal functions come from the same
parent shape, but four parameters independently control amplitude, period,
phase, and midline. Reading these parameters directly from a formula is a
powerful modeling tool used throughout science, engineering, and music.

\textbf{Transition:} ``Today we'll dissect the general form
$f(\theta)=a\sin(b(\theta+c))+d$ and learn to identify each of the four
features just by looking at the equation.''
\end{skillbox}

\vspace{0.08in}

%----------------------------------------------------------------------
% Lesson
%----------------------------------------------------------------------
\begin{skillbox}[Lesson --- Instruction Sequence]{sky}
\begin{multicols}{2}

\textbf{Step 1 --- Vertical Transformations ($a$ and $d$)} (12 min)

Begin with $y=\sin\theta$. Model $y=3\sin\theta$: every output is
multiplied by 3, so amplitude becomes 3; period and midline unchanged.
Then add $d$: $y=3\sin\theta+2$ shifts the midline to $y=2$; max is
now 5, min is $-1$. Formalize EKs 3.6.A.2 and 3.6.A.4.
Key questions: what if $a=-2$? (reflection over midline, amplitude still 2)
What if $|a|<1$? (vertical compression toward midline).

\columnbreak

\textbf{Step 2 --- Horizontal Transformations ($b$ and $c$)} (12 min)

Model $y=\sin(2\theta)$: replace $\theta$ with $2\theta$; the cycle
completes in $\pi$ instead of $2\pi$, halving the period.
Formalize EK 3.6.A.5: period $=2\pi/|b|$.
Introduce $c$: $y=\sin(\theta+\pi/4)$ shifts the graph \emph{left}
by $\pi/4$ (phase shift $=-c=-\pi/4$). Emphasize factoring:
$\sin(2\theta+\pi)=\sin(2(\theta+\pi/2))$ --- phase shift $= -\pi/2$,
not $-\pi$. Students practice factoring two or three examples.

\end{multicols}

\begin{multicols}{2}

\textbf{Step 3 --- The General Form and Summary Table} (10 min)

Write $f(\theta)=a\sin(b(\theta+c))+d$ and label all four parameters.
Complete the summary table in guided notes. Work one full example together:
$f(\theta)=-2\sin\!\bigl(\tfrac{1}{2}(\theta-\pi)\bigr)+3$ ---
amplitude $=2$, period $=4\pi$, phase shift $=\pi$ (right),
midline $y=3$, reflected ($a<0$).

\columnbreak

\textbf{Step 4 --- Sine--Cosine Equivalence} (8 min)

Establish EK 3.6.A.1: $\cos\theta=\sin(\theta+\pi/2)$.
Demonstrate: $2\cos(3\theta)=2\sin(3(\theta+\pi/6))$.
Students verify with a table of values at $\theta=0,\,\pi/6$.
Close with the key idea: \emph{any} sinusoidal function can be written
as either a sine or cosine transformation; the choice is often one of
convenience.

\end{multicols}

\smallskip
\textbf{Pacing note:} Steps 1--2 in Period 1; Steps 3--4 and group
activity in Period 2; review and exit ticket in Period 3 if needed.
\end{skillbox}

\vspace{0.08in}

%----------------------------------------------------------------------
% Active Monitoring
%----------------------------------------------------------------------
\begin{skillbox}[Active Monitoring]{redbox}
\begin{itemize}[leftmargin=1.4em, itemsep=3pt]
  \item \textbf{Phase-shift sign error (most common)}: Students read $+c$
    as the phase shift rather than $-c$. Prompt: \textit{``Rewrite
    $\sin(2\theta+\pi)$ in factored form $\sin(b(\theta+c))$. What must
    $c$ be? Then the phase shift is $-c$.''}
  \item \textbf{Period vs.\ $b$}: Students report $b$ as the period.
    Ask: \textit{``If $b=4$, the wave repeats 4 times faster than normal.
    Period $=2\pi/4=\pi/2$, not 4.''}
  \item \textbf{Amplitude and sign}: Students report a negative amplitude
    when $a<0$. Reinforce: amplitude $=|a|$ is always non-negative;
    the negative sign causes a reflection.
  \item \textbf{Midline vs.\ maximum}: Students confuse $d$ with the
    max output. Prompt: \textit{``Midline $=d$ is the \emph{average}
    of max and min. Max $= d + |a|$.''}
\end{itemize}
\end{skillbox}

\vspace{0.08in}

%----------------------------------------------------------------------
% Group Work & Differentiation
%----------------------------------------------------------------------
\begin{skillbox}[Group Work \& Differentiation]{redbox}
Students work in groups of 3--4 on the tiered activity (teacher-assigned
or student self-selected with guidance).

\begin{itemize}[leftmargin=1.4em, itemsep=4pt]
  \item \textbf{Tier R (Reinforce)}: Given five equations of the form
    $a\sin(b\theta)+d$ (no phase shift), students complete a parameter
    table (amplitude, period, midline, max, min) with sentence frames
    to guide written explanations.
  \item \textbf{Tier A (Apply)}: Given equations in factored form
    $a\sin(b(\theta+c))+d$, students identify all four parameters,
    determine the sine--cosine equivalence, and match each equation to
    a provided sinusoidal graph card (no sketching).
  \item \textbf{Tier E (Extend)}: Given equations in unfactored form
    (e.g., $-3\cos(2\theta-\pi/3)+1$), students factor the argument,
    identify all parameters, write each as a sine function, and write
    a verbal transformation description in the correct algebraic order.
\end{itemize}
\end{skillbox}

\vspace{0.08in}

%----------------------------------------------------------------------
% Individual Work & Assessment
%----------------------------------------------------------------------
\begin{skillbox}[Individual Work \& Assessment]{redbox}
Exit ticket administered independently (final 5--7 minutes of Period 2 or 3).

\begin{enumerate}[leftmargin=1.4em, itemsep=4pt]
  \item Given $f(\theta)=-3\sin\!\bigl(2(\theta+\tfrac{\pi}{4})\bigr)-1$,
    students identify amplitude, period, phase shift, and midline.
  \item Multiple-choice: select the equation matching a verbal description
    of four sinusoidal parameters.
\end{enumerate}

\textbf{Data use:} Sort exit tickets by error type --- phase-shift sign
errors, period-formula errors, amplitude-sign errors. Use distribution
to determine whether a brief re-teach opener is needed at Lesson 3.7.
\end{skillbox}

\vspace{0.08in}

%----------------------------------------------------------------------
% Reinforcement & Extension
%----------------------------------------------------------------------
\begin{skillbox}[Reinforcement \& Extension]{goldbox}
\textbf{Homework (Lesson 3.6):} 6--7 problems covering: identifying all
four parameters from sine and cosine equations in factored and unfactored
form, computing period and amplitude, identifying phase shift direction,
and 2 AP-style multiple-choice items. Optional extension: rewrite a given
cosine function as a sine function with an explicit phase shift.

\textbf{Preview of Lesson 3.7 --- Sinusoidal Function Graphs:}
Ask students: \textit{``Now that you can read amplitude, period, midline,
and phase shift from an equation, what information would you need to
produce the actual graph? In Lesson 3.7 we'll go in both directions ---
from equation to graph, and from graph to equation.''}
\end{skillbox}

\end{document}
